\newpage
\section{Appendix: Liturgical Resolution}
This is All Saints, Solemnity, white, in the national calendar for the dioceses of the United States of America, Sunday 1 November 2026. It replaces the Thirty-first Sunday in Ordinary Time. The surrounding Sunday cycle is A, while these annual readings are common to A, B and C. The weekday cycle is inapplicable. The governing books are the U.S. \work{Roman Missal}, Third Edition (2011), with the third typical Latin edition's amended state, and the U.S. second typical \work{Lectionary for Mass}, no. 667. There is no separate appointed Vigil Mass; an anticipated celebration uses this Mass.\footnote{USCCB, \work{Liturgical Calendar for the Dioceses of the United States of America 2026}, printed pp. 5,42; \work{Ordo lectionum Missae} (1981), no. 667, printed p. 313; official U.S. daily readings, 1 November 2026.}

The permanent identity is PC-R07-ABC. The canonical reusable Missal evidence belongs to the same edition's All Saints shared owner; the composition audit and instance record identify its relationship to this leaf. No local calendar or ritual overlay has been established. Required prayers and readings are distinguished from general permissions and documented local selections below; none of the latter is asserted.

\begin{branchtable}
\bid{annual-lectionary-667} & U.S. appointments and OLM 667 & Required & Readings, psalm, acclamation & One annual set; no short form or closed alternative.\\
\bid{missal-antiphons} & GIRM 48,87; All Saints & Appointed alternative & Entrance, Communion & Missal antiphons treated; actual singing choice unknown.\\
\bid{other-lawful-chants} & GIRM 48,61--62,74,87 & Permitted & Chants & Open classes; no specific local chant established.\\
\bid{preface-with-ep-i} & GIRM 365(a) & Permitted & Preface, EP & Roman Canon with the proper Preface.\\
\bid{preface-with-ep-iii} & GIRM 365(c) & Permitted & Preface, EP & EP III with the proper Preface.\\
\bid{preface-with-ep-ii} & GIRM 365(b) & Conditional & Preface, EP & Retain the proper Preface; no special circumstance or local selection established.\\
\bid{ep-iv} & GIRM 365(d) & Not applicable & EP & Invariable Preface; this solemnity has its own.\\
\bid{all-saints-solemn-blessing} & Missal, referred blessing 18 & Permitted & Conclusion & Named option; no actual selection established.\\
\bid{gloria-and-creed} & All Saints proper rubrics & Required & Ordinary & Both said; not additional proper texts.\\
\bid{sunday-sprinkling} & GIRM 51 & Permitted & Entrance rites & General Sunday possibility; no actual choice established.\\
\bid{sequence} & Inspected proper and OLM & Not applicable & Sequence & None appointed.\\
\bid{proper-ep-insert} & Inspected proper & Not applicable & EP insert & No festival-specific insert identified.\\
\bid{local-ritual-overlay} & Competent local authority needed & Conditional & Affected units & Unresolved beyond national-calendar scope.\\
\end{branchtable}

\section{Appendix: Scope and Qualifications}
The study explains these annual All Saints appointments through two whole-Mass interpretations. Each combines checked commentary on the biblical passages with original theological comparisons among the Mass's texts. No Father is attributed an exposition of this complete modern Mass. The bounded reception search and exact source roles are recorded in the canonical research; the substantive accounts here use Augustine, Aquinas, Bede and Chrysostom, with the received Victorinus commentary as supplementary evidence.

The Douay--Rheims/Challoner Scripture is historical public-domain study English in the United States, not the approved text proclaimed at this celebration. The U.S. Lectionary controls proclamation. NPNF, ANF and Edward Marshall's Bede translations retain their own historical public-domain status; the project license does not relicense their wording. The English patristic passages were checked in the retained CCEL deliveries, without a new original-language collation. Aquinas was checked in the identified electronic Latin reportatio; the two Matthew deliveries identify different printed textual bases.

This is a rights-limited study. Full protected Latin and approved English prayer and composed-antiphon bodies are omitted; exact source locators take their place. The secondary 2002 Latin reproduction and official change list support qualified evidence rather than authentication of every 2008 altar-book word. A current authoritative U.S. publisher exemplar for every oration and the blessing remains unavailable after the recorded second search. The secondary English witness has unfinished references and an older Collect conclusion. The official antiphon and Preface excerpts have narrower authority for their own texts. Original commentary on prayer subjects does not supply a substitute translation. Rights and evidence limits are retained in the canonical source records.

Chronological alternatives follow the generated canonical record and the historical sources explained in the preceding appendix. The historical encyclopedia articles are dated scholarship. The national occurrence and branches are bounded to the stated date and territory. No parish enactment, exhaustive corpus search, complete altar-book collation or external ecclesiastical approval is asserted.

\section{References}
\begin{itemize}[leftmargin=1.3em,itemsep=.4em]
\item \textbf{Scripture.} \work{Douay--Rheims Bible}, Challoner text, Project Gutenberg eBook 1581, retained verse-text edition. Appointed and contextual passages are cited at their use.
\item \textbf{Augustine.} \work{Our Lord's Sermon on the Mount}, I.1--4, §§1--12, trans. William Findlay; \work{Tractates on the First Epistle of John}, IV.4--7,9, trans. H. Browne; \work{Expositions on the Psalms}, Psalm 23, §§1--6, in NPNF, first series, vols. VI, VII and VIII, historical English in the retained CCEL deliveries. Psalm XXIV in the English heading is Latin XXIII.
\item \textbf{Thomas Aquinas.} \work{Super Psalmos} 23, nn. 2--4, Corpus Thomisticum, cps21, paragraphs 86989--86991. \work{Super Matthaeum}, reportatio of Peter of Andria, cap. 5, lectio 2, paragraph 87340, cma0500, header identifying Turin 1951; continuation, paragraphs 87564--87566, cma0510, header identifying Leonine 1969. Latin electronic deliveries inspected 26 October 2026; no print-image collation claimed. \url{https://www.corpusthomisticum.org/cps21.html}; \url{https://www.corpusthomisticum.org/cma0500.html}; \url{https://www.corpusthomisticum.org/cma0510.html}.
\item \textbf{Bede.} \work{Explanation of the Apocalypse}, trans. Edward Marshall (Oxford and London: James Parker, 1878), I, on chapter 7, especially verses 2,4,9--14, pp. 44--53. Marshall's biblical catchwords are not the appointed English.
\item \textbf{John Chrysostom.} \work{Homilies on the Gospel of St. Matthew}, XV.2--9 and XXXVIII.3, trans. George Prevost, revised M. B. Riddle, NPNF, first series, X; retained historical CCEL delivery.
\item \textbf{Received Victorinus commentary.} \work{Commentary on the Apocalypse}, seventh chapter, comments 2,9, ANF VII, trans. Robert Ernest Wallis. The edition's elucidations identify recension and attribution difficulties; the received text is cited in that qualified role.
\item \textbf{Liturgical books and controls.} \work{Roman Missal}, Third Edition, U.S. (2011), All Saints and Solemn Blessing no. 18; qualified secondary Latin 2002 reproduction, PDF 512--514 and 377--378, with the official 2008 \work{Variationes}, \work{Notitiae} 503--504, pp. 368--387. ICEL, \work{Antiphonary}, All Saints, p. 134; ICEL, \work{The glory of Jerusalem, our mother}, two-page publisher Preface. \work{Ordo lectionum Missae} (1981), no. 667, p. 313. USCCB, \work{Liturgical Calendar 2026}, pp. 5,42; \href{https://bible.usccb.org/bible/readings/110126.cfm}{daily readings for 1 November 2026}; U.S. GIRM, nos. 48,51,61--62,74,87,365, exact cached chapters inspected for this research.
\item \textbf{Historical orientation.} USCCB, NABRE, ``The Book of Psalms,'' introduction. E. Jacquier, ``Gospel of St. Matthew,'' \work{Catholic Encyclopedia} X (1911). Walter Drum, ``Epistles of Saint John,'' VIII (1910). A. Durand, ``The New Testament,'' XIV (1912). F. Van den Biesen, ``Apocalypse,'' I (1907). Exact date/place discussions are cited in the historical appendix.
\end{itemize}
